\documentclass[11pt]{article}
\usepackage[margin=1.15in]{geometry}
\usepackage{amsmath,amssymb,amsthm,mathtools}
\usepackage{xcolor}
\usepackage[colorlinks=true,linkcolor=blue!60!black,citecolor=blue!60!black,urlcolor=blue!60!black]{hyperref}

\newtheorem{theorem}{Theorem}[section]
\newtheorem{proposition}[theorem]{Proposition}
\newtheorem{lemma}[theorem]{Lemma}
\newtheorem{corollary}[theorem]{Corollary}
\theoremstyle{definition}
\newtheorem{definition}[theorem]{Definition}
\theoremstyle{remark}
\newtheorem{remark}[theorem]{Remark}

\newcommand{\C}{\mathbb{C}}
\newcommand{\R}{\mathbb{R}}
\newcommand{\Z}{\mathbb{Z}}
\newcommand{\w}{\omega}
\newcommand{\Sp}{\mathrm{Sp}}
\newcommand{\dg}{\dagger}
\newcommand{\dd}{\ddagger}
\newcommand{\WF}{\operatorname{WF}}

\title{Gauge Reduction and the Completion Problem in Fourth-Order Gravity:\\
PU Pairing, Covariant Real Forms, and the Conformal Jordan Boundary}
\author{GPT-5.6.sol\thanks{OpenAI model and principal programme author.}
\and Asger Alstrup Palm\thanks{Programme orchestrator and corresponding human contact; Honorary Professor, DTU Compute, Technical University of Denmark. \texttt{asger@area9.dk}.}}
\date{Open-repository preprint, July 2026}

\begin{document}
\maketitle

\begin{abstract}
We classify translation-invariant, quasifree, mode-local real-form
completions of free scalar-free Einstein--Weyl gravity,
$\mathcal L=\sqrt{-g}\,(c_1R+\alpha R_{\mu\nu}^2+\beta R^2)$ with
$\alpha=-3\beta$, linearized about Minkowski space, subject to the reduced
symplectic form, the spectral condition, and Poincar\'e covariance.  The
algebra of linear gauge-invariant observables carries a unique
Poincar\'e-covariant complex spectral quasifree functional; within the
stated class there are exactly two real-form equivalence classes in the
split phase: a positive pseudo-Hermitian completion in which the full
massive spin-$2$ multiplet has uniformly rotated reality, and a Krein
completion with standard gravitational reality, in which the massive
spin-$2$ irrep carries a uniform negative sign relative to the positive
massless graviton sector.  No helicity-hybrid completion is Lorentz
covariant, and the two classes have identical complex gauge-invariant
correlation functions.  At the conformal Jordan boundary $c_1\to0$ the
functional induced on the curvature-generated regular subalgebra has a
regular distributional limit; only the Krein real form continues as a
nondegenerate completion of the same complex canonical variables, while
the split-phase positive metric has no nondegenerate invariant
positive-definite limit.  The mechanism is a stratification:
diffeomorphism reduction preserves the Pais--Uhlenbeck (PU) pairing of the
fourth-order dynamics exactly in the helicity-$\pm2$ sector, while the
massive helicity-$\pm1,0$ polarizations survive as unpaired ghost
oscillators subject to a completion trilemma --- positivity, the spectral
condition, and standard reality: any two, never all three.  On the
conformal locus the transverse-traceless sector becomes a $\Box^2$ Jordan
block per polarization, the vector modes become ordinary massless ghosts,
and the scalar mode enters the kernel of the presymplectic form and is
removed by the emergent Weyl gauge symmetry, giving the conformal count of
six.  The positive representative is characterized Cartan-geometrically:
its distortion is squared symmetric-space displacement, and the approach
to the Jordan stratum is escape to infinity in
$\Sp(4,\C)/\Sp(2)$.  All finite-dimensional, variational, symplectic,
tensor-algebraic, and modewise distributional identities are independently
machine-verified.
\end{abstract}

\section*{AI authorship and computational accountability}
GPT-5.6.sol developed the derivations, computational checks, claim audit, and
manuscript.  Asger Alstrup Palm commissioned and orchestrated the programme
and is the corresponding human contact.  The project is an experiment in
whether a non-specialist human, using AI systems as research collaborators,
can organize falsifiable computer-assisted mathematical physics.  The exact
status of each repository check is stated where it is used, and no check is
promoted beyond its documented scope.  The manuscript remains subject to
expert review.

\section{Introduction}\label{sec:intro}

Fourth-order gravity has two well-developed responses to the Ostrogradsky
ghost.  In the PT-symmetric program of Bender and Mannheim
\cite{BM2008PRD,Mannheim2021}, a positive metric operator renders each
Pais--Uhlenbeck (PU) mode pair unitary on a Hilbert space; in the
Krein-space program, developed for the pure fourth-order scalar by Bateman
and Turok \cite{BT2026}, the state space is indefinite and probabilities
are controlled by ghost parity.  A companion series
\cite{companion1,companion2,companion3} analyzed the scalar case
exhaustively: the positive metric is canonically reconstructed and
variationally selected; particle content decouples from metric geometry;
and the two programs determine the \emph{same} complex spectral vacuum
functional, differing only in the real form --- the involution and
completion --- placed upon it.

Gravity poses a question the scalar theory cannot: the completion must be
chosen \emph{jointly} for all polarizations of a gauge theory, subject to
diffeomorphism reduction and Lorentz covariance.  This paper answers it
for linearized scalar-free Einstein--Weyl gravity.  The logic runs
\begin{equation*}
\text{reduction}\ \rightarrow\ \text{completion trilemma}\ \rightarrow\
\text{covariance classification}\ \rightarrow\ \text{Cartan geometry}\
\rightarrow\ \text{spectral bridge}\
\rightarrow\ \text{gauge-invariant correlator}\ \rightarrow\
\text{conformal boundary},
\end{equation*}
and the results are:

\begin{enumerate}
\item \emph{Reduction stratifies} (Section~\ref{sec:reduction}).  At fixed
spatial momentum the physical phase space is
\begin{equation}\label{eq:strat}
\Gamma_{\mathrm{phys}}(k)\ \cong\
2\,\Gamma_{\mathrm{PU}}(0,M)\ \oplus\ 2\,\Gamma_-(M)\ \oplus\ \Gamma_-(M),
\qquad M^2=c_1/\alpha:
\end{equation}
the two transverse-traceless (TT) polarizations are exact PU pairs
(massless plus massive branch, in the perfect-square-plus-Einstein form
with $\gamma=\alpha/2$); the massive helicity-$\pm1$ and $0$ modes are
\emph{unpaired} second-order ghost oscillators, because their massless
partners lie entirely in the diffeomorphism orbit.  PU pairing survives
gauge reduction exactly in helicity $\pm2$.

\item \emph{Completion trilemma} (Section~\ref{sec:trilemma}).  An
unpaired ghost oscillator admits a quarter-turn positive completion
($q\mapsto iq$, $p\mapsto-ip$, $\eta=e^{-\pi D}$, $D=\tfrac12(qp+pq)$)
whose physical adjoint flips the sign of the original variables.  The
three desiderata --- positive norm, positive energy (spectral condition),
and standard gravitational reality --- can be satisfied two at a time,
never all three.  All complex-symplectic real-form changes taking $H_-$ to
$H_+$ form a single quarter-turn coset $T_0\cdot SO(2,\C)$.

\item \emph{Covariance classification} (Section~\ref{sec:covariance}).
The five massive polarizations form one spin-$2$ irrep of the massive
little group; by Schur, any invariant Hermitian form on it is a multiple
of the identity, so its signature is uniform: \emph{no helicity-hybrid
completion is covariant}.  An induced-representation lemma lifts the
rest-frame statement to the Poincar\'e group: translation invariance
superselects the mass shells, covariant Hermitian structures correspond
to little-group invariant fiber forms, and the equivariant antilinear
involutions on the real-type spin-$2$ irrep fall into exactly two classes.
The total quarter-turn generator assembles into an equivariant bundle map,
so a covariant uniform-positive completion exists; and the quarter-turn on
the ghost normal mode of the TT PU pair is itself an admissible positive
diagonalizer, hence lies in the metric family of \cite{companion1} and, by
orbit constancy \cite{companion3}, defines the same vacuum functional as
the Bender--Mannheim metric: the TT and lower-helicity structures
assemble.

\item \emph{Cartan geometry of the positive representative}
(Section~\ref{sec:cartan}).  For one PU block the minimum-distortion
functional is squared Cartan displacement, $F(S)=2\|\mu(S)\|^2$ with
$\mu:\Sp(4,\C)\to\overline{\mathfrak a_+}$ the Cartan projection, and the
canonical positive diagonalizer sits on the $C_2$ Weyl-chamber wall,
$\mu(S_+)=(r,r)$ \cite{companion2}.  This is not a direct instance of
standard Mostow/Kempf--Ness theory --- the stabilizer $SO(2,\C)^2$ is not
Cartan-compatible ($H\cap K\cong\Z_2\times\Z_2$) --- and the proof runs
through the $\theta$-compatible hull $SL(2,\C)\times SL(2,\C)$.  The
Cartan picture unifies the companion imports: split orbit $\to$ Cartan
displacement $\to$ escape to infinity $\to$ Jordan boundary.

\item \emph{Spectral bridge} (Section~\ref{sec:kernel}).  The reduced
spectral kernel, with residues fixed by the reduced symplectic form,
reassembles --- modulo pure-gauge bi-distributions --- into the covariant
massive spin-$2$ projector with the single overall coefficient
$4/c_1$ and uniform ghost sign; the massless shell couples only to TT.
The complex kernel is identical in both real forms.  Together with item 3
this closes the bridge at the gravitational level: \emph{the two covariant
completions induce the same vacuum functional on the gauge-invariant
observable algebra}.

\item \emph{Gauge-invariant correlator} (Section~\ref{sec:weyl}).  The
linearized Weyl correlator is gauge independent, and the Weyl map
annihilates every $p_\mu p_\nu/M^2$ term of the massive projector: the
curvature kernel is $M$-regular and covariant, with Hadamard wavefront
directions; the fourth-order logarithm of the potential kernel appears in
the curvature correlator as its differentiated descendant.

\item \emph{Conformal Jordan boundary} (Section~\ref{sec:conformal}).  As
$c_1\to0$ at fixed $\alpha$: the TT kernel converges distributionally to
$1/\gamma$ times the $\Box^2$ Jordan kernel; the vector modes have a
smooth massless-ghost limit with fixed commutator normalization; the
scalar direction enters the kernel of the presymplectic form and becomes
Weyl-gauge (the transformation $\delta h_{\mu\nu}=2\sigma\eta_{\mu\nu}$
is a symmetry exactly at $c_1=0$), giving the conformal count $4+2+0=6$.
Because the observable algebra changes rank at the boundary, convergence
is formulated through explicit comparison maps on the curvature-generated
regular subalgebra.  The split rapidity of the TT normal-mode
decomposition underlying the positive completion diverges exactly, for
every fixed $k>0$:
$e^{r}=(\sqrt{k^2+M^2}+k)^2/M^2\sim4k^2/M^2$ --- Cartan-geometrically, the
positive representative escapes to infinity in the symmetric space.  The
positive metric has no nondegenerate invariant positive-definite limit;
the Krein real form and the spectral functional continue.
\end{enumerate}

The classification is summarized by the master theorem of
Section~\ref{sec:master}.  Throughout, gauge (Faddeev--Popov/BRST) ghosts
--- which encode the diffeomorphism quotient and are handled here by
explicit symplectic reduction --- are sharply distinguished from the
higher-derivative spin-$2$ ghost branch, which concerns the completion of
the physical excitations.

\paragraph{Verification.}
All finite-dimensional, variational, symplectic, tensor-algebraic, and
modewise distributional identities in this paper are independently
machine-verified: the second variation and sector reduction (checks
G1--G7), the trilemma and real-form classification (G8), the covariance
computations (G9), the spectral kernel and covariant reassembly including
the absolute normalization (G10), the Weyl-correlator cancellations and
symbol nonvanishing (G11), and the sectorwise conformal limits including
the exact divergence law of the split decomposition (G12).
Appendix~\ref{app:checks} maps each check to the theorem it supports;
the artifact is the root of the \texttt{area9innovation/weyl-gravity}
repository accompanying the companion series
\cite{companion1,companion2,companion3} (scripts
\texttt{symbolic/verify\_gravity\_\{reduction,completion,spectral\}.py},
run with \texttt{python3}; SymPy 1.14.0, NumPy 2.4.4, Python 3.14; cite
the repository commit hash to fix the artifact).  The
representation-theoretic and microlocal conclusions --- existence of the
Hilbert/Krein completions, Poincar\'e covariance as a representation
statement, and wavefront-set equality --- follow from the analytic
arguments given in the text and the companion papers, not from the
symbolic engine.

\paragraph{Conventions.}
Signature $(+,-,-,-)$; curvature conventions fixed so that the healthy
Einstein normalization is $c_1=-1$ (checked on the graviton kinetic term);
$\alpha=-3\beta$ (scalar-free tuning); $M^2=c_1/\alpha>0$ for $\alpha<0$;
modes $h_{\mu\nu}\sim A(t)\cos kz+B(t)\sin kz$ at spatial momentum $k$
along $z$; $\w_1=\sqrt{k^2+M^2}$, $\w_2=k$.  PU data follow
\cite{companion1,companion2}: a PU pair with parameter $\gamma$ has
Lagrangian $\tfrac{\gamma}{2}(\ddot z^2-(\w_1^2+\w_2^2)\dot
z^2+\w_1^2\w_2^2z^2)$ and commutator normalization $E'''(0)=1/\gamma$.

\section{Reduction: the helicity stratification}\label{sec:reduction}

The second variation of
$\mathcal L=\sqrt{-g}(c_1R+\alpha R_{\mu\nu}^2+\beta R^2)$ about flat
space is computed exactly, sector by sector under the rotation group about
$k$; diffeomorphism invariance of each sector Lagrangian is verified (the
gauge variation is a total derivative).

\begin{theorem}[TT sector: exact PU blocks]\label{thm:tt}
For each of the two TT polarizations the mode Lagrangian is
\begin{equation}\label{eq:LTT}
L_{\mathrm{TT}}
=\frac{\alpha}{4}\big(\ddot A+k^2A\big)^2
-\frac{c_1}{4}\big(\dot A^2-k^2A^2\big)+\frac{d}{dt}(\cdots),
\end{equation}
a perfect square plus Einstein term, with Euler--Lagrange equation
$(\partial_t^2+k^2)(\partial_t^2+k^2+M^2)A=0$, $M^2=c_1/\alpha$: an exact
PU block with mass pair $(M,0)$ and normalization $\gamma=\alpha/2$.  The
sign $\gamma<0$ (for $c_1=-1$, $\alpha<0$) is the perfect-square
convention of \cite{BT2026} and renders the \emph{massless} branch healthy
and the massive spin-$2$ branch the ghost, as required by the flat-space
spectrum of quadratic gravity \cite{Stelle}.
\end{theorem}

\begin{theorem}[Lower helicities: PU pairing broken by gauge]
\label{thm:lower}
(i) Vector sector: with the gauge-invariant $w=k\,h_{tx}+\partial_t
h_{xz}$ (one per transverse direction),
\begin{equation}
L_V=\frac{\alpha}{4}\big(\dot w^2-k^2w^2\big)-\frac{c_1}{4}w^2
+\frac{d}{dt}(\cdots):
\end{equation}
a single second-order massive mode, $\ddot w+(k^2+M^2)w=0$, with
ghost-signed kinetic normalization $\mu_V=\alpha/2$; the massless PU
partner is pure diffeomorphism gauge and is removed by the quotient.
(ii) Scalar sector: $h_{tt}$ is auxiliary; at $\alpha=-3\beta$ the reduced
dynamics is the single massive mode $\ddot p+(k^2+M^2)p=0$ with effective
kinetic normalization $\mu_S=3c_1/2$ (ghost-signed, $\beta$-independent);
for generic couplings the sector is fourth order with scalaron mass
$m_0^2=-c_1/(2(\alpha+3\beta))$, decoupling exactly at the scalar-free
tuning.  (iii) Mode count: $2\times2+2+1=7=2+5$, and the stratification
\eqref{eq:strat} holds: PU pairing survives reduction exactly in helicity
$\pm2$.
\end{theorem}

\begin{definition}[Invariant observable algebra]\label{def:Ainv}
$\mathfrak A_{\mathrm{inv}}$ denotes the algebra of linear gauge-invariant
observables of the linearized theory --- equivalently, the field algebra
of the reduced phase space \eqref{eq:strat} --- with, at $c_1=0$, the
additional quotient by the kernel of the presymplectic form (the emergent
Weyl-null direction of Section~\ref{sec:conformal}).  All uniqueness and
continuation statements in this paper are statements about states and
kernels on $\mathfrak A_{\mathrm{inv}}$; potential-field kernels
$W^+_h$ are determined only modulo pure-gauge bi-distributions and are
used as computational representatives.
\end{definition}

\begin{remark}[Polarization-enhanced stabilizer]
For the doubled TT block the normal-form stabilizer Lie algebra jumps from
$\dim_\R 4$ (one PU pair; $so(2,\C)^2$ \cite{companion1}) to $\dim_\R16$
(computed by null-space methods): frequency degeneracy across
polarizations enlarges the metric ambiguity.  $SO(2)$ helicity covariance
and Schur reduce the covariant TT metric to $S_+\otimes I_2$, the unique
helicity-covariant minimum-distortion positive PU diagonalizer
\cite{companion2}.
\end{remark}

\section{The completion trilemma}\label{sec:trilemma}

We first fix the category in which "exactly two completions" is a
theorem rather than a slogan.

\begin{definition}[Complex dynamical algebra, real form, completion]
\label{def:realform}
(i) The \emph{complex dynamical algebra} of the reduced theory is the
triple $(\mathcal V_\C,\sigma,U)$: the complexified reduced phase space
\eqref{eq:strat}, the complexified symplectic form $\sigma$, and the
Poincar\'e action $U$.  No conjugation is selected at this stage.
(ii) A \emph{real form} is an antilinear involution
$\kappa:\mathcal V_\C\to\mathcal V_\C$, $\kappa^2=1$, compatible with
$\sigma$, with the dynamics, and with $U$.
(iii) A \emph{completion} is a Hilbert or Krein topology on the
$\kappa$-fixed-point real subspace, together with the corresponding
adjoint and quasifree (GNS/Fock) construction.
(iv) Two completions are \emph{equivalent} if they are related by a
$U$-equivariant complex-symplectic automorphism preserving the spectral
covariance --- in particular, by a symplectic stabilizer commuting with
the dynamics.  Throughout, "mode-local" means block-diagonal over the
helicity--momentum decomposition of \eqref{eq:strat}.
\end{definition}

\begin{remark}[States vs.\ covariances]\label{rem:states}
Positivity of a state presupposes an involution.  The two constructions
classified below therefore do \emph{not} define the same positive state
on one fixed $*$-algebra; they define the same complex bilinear spectral
covariance on the complex dynamical algebra, turned into states by two
different real forms with two different positivity notions.  All "same
functional" statements in this paper are statements of the second kind.
\end{remark}

Each unpaired ghost oscillator, $H_-=-\tfrac12(p^2+\Omega^2q^2)$, poses
the completion problem in miniature.

\begin{theorem}[Unpaired-ghost completion]\label{thm:trilemma}
Let $D=\tfrac12(qp+pq)$ and $\rho=e^{-\frac{\pi}{2}D}$.  Then:
(i) $\rho q\rho^{-1}=iq$, $\rho p\rho^{-1}=-ip$, and
$\rho H_-\rho^{-1}=+\tfrac12(p^2+\Omega^2q^2)$;
$\eta=\rho^{\dg}\rho=e^{-\pi D}$ is formally positive.
(ii) The physical adjoint is $q^{\dd}=\eta^{-1}q^{\dg}\eta=-q$,
$p^{\dd}=-p$: the original amplitude is $\dd$-anti-Hermitian; $iq$ is the
observable.  (iii) Trilemma: positive norm, positive energy, and standard
reality can be realized two at a time, never all three ---
\begin{center}
\begin{tabular}{c|c|c|c}
completion & positive norm & positive energy & standard reality\\
\hline
ordinary Hilbert & \checkmark & $\times$ & \checkmark\\
Krein & $\times$ & \checkmark & \checkmark\\
quarter-turn (PT contour) & \checkmark & \checkmark & $\times$
\end{tabular}
\end{center}
(for the first row, $p^2+\Omega^2q^2$ built from a self-adjoint canonical
pair has positive unbounded spectrum, so $H_-$ is unbounded below).
(iv) Classification: all complex-symplectic real-form changes taking $H_-$
to $H_+$ form the single coset
$\{T\in\Sp(2,\C):TA_+T^{-1}=-A_+\}=T_0\cdot SO(2,\C)$,
$T_0=\operatorname{diag}(i,-i)$.
\end{theorem}

An unpaired gravitational ghost therefore cannot simultaneously satisfy
positivity, the spectral condition, and the standard real structure.  The
choice among the three completions is exactly a choice of real form.

\section{Covariance classification}\label{sec:covariance}

Sectorwise choices are constrained globally: the massive helicities
$(\pm2,\pm1,0)$ form one irreducible spin-$2$ multiplet of the massive
little group.  The following lemma lifts rest-frame statements to the
Poincar\'e group; it is the analytic complement to the finite-dimensional
computations below.

\begin{lemma}[Induced-representation structure]\label{lem:wigner}
(i) The massive one-particle space is the Wigner representation induced
from the spin-$2$ irrep of $SO(3)$ over the mass shell
$p^2=M^2$, $p^0>0$; the massless TT space is induced from the
helicity-$\pm2$ representations over $p^2=0$.
(ii) Translation invariance of a quasifree two-point functional forces
its momentum-space kernel to be supported on the joint spectrum; since
the shells $p^2=M^2$ and $p^2=0$ are disjoint $U$-orbits, no
translation-invariant Hermitian structure mixes them: the completion
problem decomposes over shells.
(iii) A Poincar\'e-covariant Hermitian structure on an induced
representation is equivalent to a little-group invariant Hermitian form
on the inducing fiber (boosts act by fiber transport, so invariance under
the full group reduces to invariance at one point of the orbit).
(iv) The spin-$2$ irrep of $SO(3)$ is of real type; its equivariant
antilinear involutions form exactly two classes, $\kappa$ and
$T_0\kappa$ (standard and quarter-turn-rotated reality), because any two
equivariant involutions differ by an equivariant linear map, which by
Schur is a scalar, and antilinearity fixes the scalar to $\pm1$ up to
phase redefinition.
(v) The momentum-dependent canonical pairs $(q_a(p),p_a(p))$ of the five
massive polarizations assemble into an equivariant bundle map over the
shell: $D_{\mathrm{tot}}(p)=\sum_a\tfrac12(q_ap_a+p_aq_a)(p)$ transforms
by fiber transport, so rest-frame $SO(3)$ invariance of
$D_{\mathrm{tot}}$ implies invariance under the boosts as well, and the
resulting real-form change commutes with the full Poincar\'e action.
\end{lemma}

\begin{proof}[Proof sketch]
(ii) is the spectral-support argument for quasifree translation-invariant
kernels; (iii) is the standard unitarity criterion for induced
representations; (iv) is Schur plus the real-type classification of
$SO(3)$ irreps (integer spin); (v) follows because the polarization
bundle over either shell is the associated bundle of the little-group
principal bundle and $D_{\mathrm{tot}}$ is built $SO(3)$-equivariantly
from its fiber data.  The fiber statements consumed by (iii)--(v) are the
machine-checked commutant and invariance computations cited in
Theorems~\ref{thm:nohybrid}--\ref{thm:tworeal}.
\end{proof}

\begin{theorem}[No covariant hybrid]\label{thm:nohybrid}
The space of $SO(3)$-invariant Hermitian forms on the five-dimensional
spin-$2$ irrep is one-dimensional, $\eta_{\mathrm{mass}}=cI_5$ (computed
by solving the commutant equations).  Hence, by
Lemma~\ref{lem:wigner}(iii), any Poincar\'e-covariant Hermitian form on
the massive one-particle space has uniform signature, and the
helicity-hybrid assignment $(+,+,-,-,-)$ --- TT-positive with
lower-helicity Krein --- violates invariance explicitly and is not
Lorentz covariant.
\end{theorem}

\begin{theorem}[Two covariant real forms]\label{thm:tworeal}
(i) The total quarter-turn generator
$D_{\mathrm{tot}}=\sum_{a=1}^{5}\tfrac12(q_ap_a+p_aq_a)$ over the massive
polarization oscillators is $SO(3)$-invariant and, by
Lemma~\ref{lem:wigner}(v), assembles into a Poincar\'e-equivariant bundle
map over the massive shell; hence
$\eta_{\mathrm{mass}}=e^{-\pi D_{\mathrm{tot}}}$ defines a covariant
uniform-positive pseudo-Hermitian completion with uniformly rotated
massive reality (the massive multiplet is $\dd$-anti-Hermitian; $i$ times
the multiplet is observable), realizing the rotated involution class of
Lemma~\ref{lem:wigner}(iv).  (ii) The Krein completion with standard
gravitational reality is covariant, and by Theorem~\ref{thm:nohybrid} its
indefiniteness lives \emph{between} irreps, never inside the massive
multiplet: the invariant form on the massive spin-$2$ one-particle irrep
is definite, and the Krein structure is
\begin{equation}
\mathcal H_{\mathrm{1p}}
=\mathcal H_{0,+2}\ \oplus\ \mathcal H_{M,-5},
\qquad\text{signature } (+,+;\,-,-,-,-,-),
\end{equation}
the direct sum of the positive massless graviton sector and the massive
spin-$2$ irrep carrying uniform negative sign.  (iii) Assembly: the
quarter-turn applied to the ghost normal
mode of the TT PU pair is an admissible positive diagonalizer
($T'^{\top}JT'=J$ and $T'^{\top}G_{\mathrm{PU}}T'$ real positive
definite), hence lies in the solution family $S_+\cdot\mathrm{Stab}$ of
\cite{companion1}; by orbit constancy \cite{companion3} it defines the
same vacuum functional as the Bender--Mannheim metric.  The TT and
lower-helicity structures of the positive class are mutually consistent.
\end{theorem}

Stabilizer choices (the enlarged TT stabilizer of
Section~\ref{sec:reduction}, the $SO(2,\C)$ freedom of
Theorem~\ref{thm:trilemma}(iv), the $W$-freedom of \cite{companion1})
produce representatives \emph{within} these classes, not additional
physical classes: by orbit constancy the vacuum ray, and hence the
quasifree functional, is unchanged.

\begin{remark}[Field-level meaning of the positive class]\label{rem:field}
At field level the positive completion is constructed in the order:
(1) modewise complex-symplectic real-form change
(Theorem~\ref{thm:trilemma}); (2) positive one-particle Hermitian
structure on the polarization spaces (Theorems~\ref{thm:nohybrid},
\ref{thm:tworeal}); (3) the corresponding quasifree (GNS)
representation.  The metric-operator notation
$\eta=e^{-\pi D_{\mathrm{tot}}}$ is formal shorthand valid on a common
algebraic domain; it is \emph{not} asserted to be a globally defined
invertible operator on the na\"ive Fock space --- indeed, by the
identity-embedding obstruction theorems of \cite{companion2,companion3},
the positive quasifree representation of the field theory is disjoint
from the na\"ive Fock representation \emph{under the identity embedding
of the auxiliary canonical variables} (the Dyson-transported physical
completions are pointed-unitarily equivalent), and the representation is
defined by step (3), not by
conjugating inside a pre-existing Hilbert space.
\end{remark}

\section{Cartan geometry of the positive representative}\label{sec:cartan}

The positive class has a canonical representative, and its selection is
geometric, not merely variational.  Recall from \cite{companion2} the
minimum-distortion functional $F(S)=\|\log S^{\dg}S\|_F^2$ on the coset
of positive PU diagonalizers.

\begin{proposition}[Cartan-projection form of the distortion]
\label{prop:cartan}
For one PU block, $F(S)=2\|\mu(S)\|^2$, where
$\mu:\Sp(4,\C)\to\overline{\mathfrak a_+}$ is the Cartan projection and
$\|\cdot\|$ the Weyl-invariant Euclidean norm on the chamber.  The
canonical positive diagonalizer satisfies
\begin{equation}
\mu(S_+)=(r,r),
\end{equation}
a point on the $C_2$ Weyl-chamber wall: the distortion of $S_+$ is
literally squared symmetric-space displacement in $\Sp(4,\C)/\Sp(2)$,
and $S_+$ is the point of the diagonalizer coset closest to the origin.
\end{proposition}

\begin{remark}[Not a standard Kempf--Ness instance]\label{rem:notKN}
The minimization is \emph{not} a direct application of Mostow or
Kempf--Ness theory: the stabilizer $H\cong SO(2,\C)^2$ of the normal-form
data is not compatible with the canonical Cartan involution $\theta$ ---
its intersection with the maximal compact subgroup is
$H\cap K\cong\Z_2\times\Z_2$, not a maximal compact subgroup of $H$.  The
proof in \cite{companion2} instead passes through the $\theta$-compatible
hull $C\cong SL(2,\C)\times SL(2,\C)\supset H$, whose orbit is totally
geodesic ($\mathbb H^3\times\mathbb H^3$), and projects orthogonally in
Cartan norm.  This twist is intrinsic to the PU normal form, not an
artifact of coordinates.
\end{remark}

The Cartan picture does three jobs in this paper.  First, it explains
canonical representative selection: $S_+$ is geometrically distinguished
inside the positive coset, while orbit constancy \cite{companion3}
guarantees that this selection is a choice of representative, not a third
physical class.  Second, it converts the equal-frequency degeneration
into geometry: the rapidity divergence of
Theorem~\ref{thm:term}, $r\sim\log(4k^2/M^2)$, is the statement
$\|\mu(S_+)\|\to\infty$ --- \emph{the positive representative escapes to
infinity in the noncompact symmetric space as the dynamics approaches the
Jordan stratum}.  Third, it unifies the companion imports used here:
minimum distortion, orbit constancy, and the Jordan obstruction are the
value, the fibers, and the boundary behaviour of one map $\mu$ restricted
to the split orbit.

\section{The reduced spectral kernel and the bridge}\label{sec:kernel}

\begin{theorem}[Helicity kernel with symplectic residues]\label{thm:kernel}
The reduced spectral kernel (spectral-support-positive: supported in
$p^0>0$; its residues are of both signs in the ghost sectors, so no
distributional positivity is implied) is
\begin{align}
W^+_{\mathrm{TT}}(t)
&=\frac{1}{\gamma M^2}
\left[\frac{e^{-i\w_1t}}{2\w_1}-\frac{e^{-i\w_2t}}{2\w_2}\right],
\qquad \gamma=\frac{\alpha}{2},\\
W^+_{V}(t)=\frac{e^{-i\Omega t}}{2\mu_V\Omega},
\quad
W^+_{S}(t)&=\frac{e^{-i\Omega t}}{2\mu_S\Omega},
\qquad
\Omega=\sqrt{k^2+M^2},\quad \mu_V=\frac{\alpha}{2},\ \mu_S=\frac{3c_1}{2},
\end{align}
each a bisolution with the commutator fixed by the reduced symplectic
form ($E'''(0)=1/\gamma$ for TT; $E'(0)=-1/\mu$ for the second-order
modes; the residues are \emph{derived}, not assumed from propagator
guesses).  Note that on the tuning $\alpha M^2=c_1$ the TT prefactor is
$M$-independent: $1/(\gamma M^2)=2/(\alpha M^2)=2/c_1$.
\end{theorem}

\begin{theorem}[Covariant reassembly]\label{thm:reassembly}
Let $\Pi^{(2,M)}_{\mu\nu\rho\sigma}
=\tfrac12(P_{\mu\rho}P_{\nu\sigma}+P_{\mu\sigma}P_{\nu\rho})
-\tfrac13P_{\mu\nu}P_{\rho\sigma}$, $P_{\mu\nu}=\eta_{\mu\nu}-p_\mu
p_\nu/M^2$, on the massive shell.  The gauge-invariant contractions are
\begin{equation}
\Pi_{xy,xy}=\tfrac12,\qquad
c\,\Pi\,\bar c\big|_{w}=\frac{(E^2-k^2)^2}{2M^2}=\frac{M^2}{2},\qquad
\Pi\big|_{(h_{xx}+h_{yy})/2}=\tfrac16 ,
\end{equation}
where the complex-basis vector invariant is
$\tilde w=-ik\,h_{tx}-iE\,h_{xz}$.  These match the reduced residue ratios
$1:M^2:\tfrac13$ of Theorem~\ref{thm:kernel} exactly, and the absolute
normalizations fix a \emph{single} overall covariant coefficient,
\begin{equation}
C\cdot\tfrac12=\frac{1}{\gamma M^2},\qquad
C\cdot\frac{M^2}{2}=\frac{1}{\mu_V},\qquad
C\cdot\tfrac16=\frac{1}{\mu_S}
\quad\Longrightarrow\quad
C=\frac{4}{c_1},
\end{equation}
all three exactly on the tuning $c_1=\alpha M^2$, with uniform ghost
sign: the five massive residues assemble into one $SO(3)$-invariant
projector.  Consequently the potential kernel, as an equivalence class
\begin{equation}
[\widetilde W^+_h]\ \in\
\mathcal D'(\mathrm{Sym}^2\otimes\mathrm{Sym}^2)\big/
\{\text{pure-gauge bi-distributions}\},
\end{equation}
has the exact representative
\begin{equation}\label{eq:exactkernel}
\widetilde W^+_{\mu\nu\rho\sigma}(p)
=\frac{4}{c_1}\,\theta(p^0)\left[
\Pi^{(2,M)}_{\mu\nu\rho\sigma}(p)\,\delta(p^2-M^2)
-\Pi^{(2,0)}_{\mu\nu\rho\sigma}(p;n)\,\delta(p^2)
\right],
\end{equation}
where $\Pi^{(2,0)}(p;n)$ is the massless spin-$2$ projector in any
reference frame $n$; the class $[\widetilde W^+_h]$ is
frame-independent because changing $n$ changes $\Pi^{(2,0)}$ by a
pure-gauge kernel.  The massless shell couples only to TT (the vector and
scalar invariants contract to zero with the massless polarizations).
Since $4/c_1=(4/\alpha)/M^2$, the coefficient may equivalently be written
$\mathcal N/M^2$ with $\mathcal N=4/\alpha$; written either way it
behaves as $M^{-2}[\delta(p^2-M^2)-\delta(p^2)]\to-\delta'(p^2)$ in the
conformal limit --- one power of $M^{-2}$, exactly what the Jordan limit
of Section~\ref{sec:conformal} requires.
Equivalently and frame-freely: the exact curvature kernel
$\widetilde{\mathcal W}^+_{CC}=\mathcal D_C\widetilde W^+_h\mathcal
D_C^{\top}$ of Section~\ref{sec:weyl} contains no gauge terms and no
singular projector pieces, and determines the state on
$\mathfrak A_{\mathrm{inv}}$.
\end{theorem}

\begin{theorem}[Real-form independence; the gravitational bridge]
\label{thm:bridge4}
The complex spectral kernel is identical in the two covariant real forms:
for each unpaired ghost the quarter-turn completion's physical Wightman
function equals the Krein spectral value,
$\langle(i\tilde w)(t)\,(i\tilde w)(0)\rangle
=e^{-i\Omega t}/(2\mu\Omega)$, and for the TT pairs the statement is the
orbit constancy and spectral identification of
\cite{companion3}.  Hence the positive pseudo-Hermitian and Krein
completions of scalar-free Einstein--Weyl gravity induce the same complex
quasifree functional on the gauge-invariant observable algebra
$\mathfrak A_{\mathrm{inv}}$; they differ in involution and completion
only.
\end{theorem}

\section{The gauge-invariant Weyl correlator}\label{sec:weyl}

Curvature observables remove both gauge dependence and projector
singularities.

\begin{theorem}[Weyl correlator]\label{thm:weyl}
Let $C^{(1)}[h]$ be the linearized Weyl tensor (traceless; validated
symbolically) and $\mathcal
W^+=\mathcal D_C\,W^+_h\,\mathcal D_C^{\top}$ the induced curvature
kernel.  Then: (i) the linearized Riemann map annihilates pure gauge
$h=p\otimes\xi+\xi\otimes p$ identically, so $\mathcal W^+$ is gauge
independent; (ii) the full Weyl--Weyl contraction of $\Pi^{(2,M)}$ equals
that of $\Pi_0$ (all $P\to\eta$): every $p_\mu p_\nu/M^2$ and $/M^4$ term
of the massive projector is annihilated --- the curvature kernel is
$M$-regular, and the five massive helicities enter only through the
$SO(3)$-invariant flat projector; (iii) both real forms give the same
complex curvature functional (Theorem~\ref{thm:bridge4}); (iv)
$\WF(\mathcal W^+)=\mathcal C^+$, the standard positive-frequency Hadamard
set with tensor polarization constraints.  The fourth-order logarithmic
singularity belongs to the potential kernel, $W_h^+\sim\log\rho$, while
$\mathcal W^+_{CC}\sim\partial^4\log\rho$ is its differentiated
descendant.  Contact terms (momentum-space polynomials) are separated from
the nonlocal two-point distribution.
\end{theorem}

\begin{proof}[Proof of (iv)]
Since $\mathcal W^+=\mathcal D_C W^+_h\mathcal D_C^\top$ with $\mathcal
D_C$ a differential operator, $\WF(\mathcal W^+)\subseteq\WF(W^+_h)
=\mathcal C^+$ \cite{companion3}.  The reverse inclusion does not follow
from symbol nonvanishing alone: for tensor-valued distributions,
polarization-dependent cancellations are possible, and the
massive-minus-massless divided difference in \eqref{eq:exactkernel} does
cancel the leading ($\rho^{-1}$-type) scalar Hadamard singularity.  Three
observations close the gap.  First, the cancellation \emph{lowers the
singular order} --- the divided difference replaces $\delta(p^2)$-type
decay by $\delta'(p^2)$-type decay in overall strength, i.e.\ the
$\log\rho$ term of the fourth-order parametrix --- but removes no
covector directions: both shells project to the same characteristic cone
$\mathcal C^+$, and the difference of two distributions singular on the
same cone remains singular there unless the full symbol cancels, which
the explicit Hadamard expansion of \cite{companion3} excludes: the
logarithmic coefficient is the universal $1/(8\pi^2)$ and is nonzero.
Second, in the language of Dencker's polarization sets, the polarization
set of $W^+_h$ over each null covector is the physical spin-$2$
polarization subspace, and the linearized-Weyl principal symbol
$\sigma_C(p)$ is injective on it: on the massive shell
$\sigma_C\Pi^{(2,M)}\sigma_C^{\top}$ has the nonzero full contraction of
part (ii), and on the massless shell $\sigma_C$ is nonzero on TT
polarizations while annihilating null pure gauge (both machine-checked,
G11c/G11e).  Third, applying $\sigma_C$ therefore maps the polarization
set injectively rather than to zero, so every direction of
$\mathcal C^+$ carried by $W^+_h$ on physical polarizations survives in
$\mathcal W^+$.  Hence $\WF(\mathcal W^+)=\mathcal C^+$.
\end{proof}

\section{The conformal Jordan boundary}\label{sec:conformal}

Take $c_1\to0$ at fixed $\alpha$ (so $M^2\to0$) \emph{before} imposing any
completion.

\begin{theorem}[Sectorwise conformal limit]\label{thm:conformal}
For every fixed spatial momentum $k>0$:
(i) TT: distributionally,
$[\delta(p^2-M^2)-\delta(p^2)]/M^2\to-\delta'(p^2)$; per polarization the
physical kernel converges to the $\Box^2$ Jordan kernel
\begin{equation}
W^+_{\mathrm{TT},0}(t)=-\frac{(1+ikt)\,e^{-ikt}}{4\gamma k^3},
\qquad \gamma=\frac{\alpha}{2},
\end{equation}
(four TT configuration modes; the $1/\gamma$ prefactor is mandatory at
fixed $\alpha$ --- the unit-normalized kernel $-(1+ikt)e^{-ikt}/(4k^3)$
is the divided difference before the symplectic normalization is
applied).
(ii) Vector: $\delta(p^2-M^2)\to\delta(p^2)$ smoothly with \emph{fixed}
commutator normalization $\mu_V=\alpha/2$: ordinary second-order massless
ghost modes, not Jordan partners.
(iii) Scalar: the coefficient of $\dot p^2$ in the reduced Lagrangian,
$\mu_S/2=3c_1/4$, tends to zero, and the Weyl transformation
$\delta_\sigma h_{\mu\nu}=2\sigma\eta_{\mu\nu}$ is a gauge symmetry of
the scalar-free action exactly at $c_1=0$ (and not for $c_1\neq0$): the
scalar direction enters the kernel of the presymplectic form and is
removed by the emergent Weyl gauge symmetry.
The conformal count is $4+2+0=6$.
\end{theorem}

\begin{remark}[Zero mode]
At $k=0$ the massless branch has zero frequency and the normal-mode
analysis degenerates for a different, infrared reason.  On $\R^3$ the
$k=0$ mode is excluded from the one-particle structure by the
test-function framework (smearing with $\tilde f(0)=0$ where needed, as
in the infrared treatment of \cite{companion3}); on compact spatial
sections it would require separate treatment, which we do not attempt
here.
\end{remark}

\begin{definition}[Comparison maps at the boundary]\label{def:compmaps}
The algebras $\mathfrak A_{c_1\neq0}$ and $\mathfrak A_0
=\mathfrak A_{\mathrm{pre}}/\mathfrak I_{\mathrm{Weyl}}$ are not equal,
so "$\omega_{c_1}\to\omega_0$" requires comparison maps.  Let
$\mathcal O_{\mathrm{reg}}$ be the \emph{curvature-generated algebra}:
the algebra generated by the smeared linearized Weyl tensor.  For every
$c_1$ (including $0$) there is a canonical map
$\iota_{c_1}:\mathcal O_{\mathrm{reg}}\to\mathfrak A_{c_1}$, because the
Weyl observable annihilates both diffeomorphism gauge
(Theorem~\ref{thm:weyl}(i)) and, at $c_1=0$, the emergent Weyl-null
direction ($\iota_0$ factors through the quotient by
$\mathfrak I_{\mathrm{Weyl}}$).  The boundary statement of
Theorem~\ref{thm:master} is
\begin{equation}\label{eq:convergence}
\omega_{c_1}\circ\iota_{c_1}\ \longrightarrow\ \omega_0\circ\iota_0
\qquad\text{distributionally on }\mathcal O_{\mathrm{reg}},
\end{equation}
and \emph{not} convergence on the complete split-phase observable
algebra.
\end{definition}

\begin{remark}[Regularity holds on the quotient, not before it]
\label{rem:quotient}
The raw scalar kernel $W^+_S=e^{-i\Omega t}/(2\mu_S\Omega)$ with
$\mu_S=3c_1/2$ diverges like $1/c_1$: the full reduced-coordinate kernel
of the split phase has \emph{no} regular limit.  What converges is the
family of presymplectic phase spaces
\begin{equation}
(\Gamma_{c_1},\Omega_{c_1})\ \longrightarrow\
(\Gamma_0/\ker\Omega_0,\ \bar\Omega_0),
\end{equation}
together with the pulled-back functionals \eqref{eq:convergence} on the
curvature-generated algebra: the divergent direction is exactly the one
that enters $\ker\Omega_0$, is annihilated by $\iota$-observables, and is
quotiented by the emergent Weyl symmetry.  The observable algebra changes
rank at the boundary; this is a structural feature of the conformal
locus, not an artifact.
\end{remark}

\begin{center}
\begin{tabular}{c|c|c|c}
 & $\pm2$ & $\pm1$ & $0$\\
\hline
$c_1\neq0$ & PU pair & massive ghost & massive ghost\\
$c_1=0$ & $\Box^2$ Jordan & massless ghost & Weyl gauge/null
\end{tabular}
\end{center}

\begin{theorem}[Termination of the positive real form]\label{thm:term}
The uniform positive completion is built on the split normal-mode
decomposition of the TT PU pairs, whose rapidity is
\begin{equation}
r(k,M)=\log\frac{(\w_1+\w_2)^2}{\w_1^2-\w_2^2}
=\log\frac{\big(\sqrt{k^2+M^2}+k\big)^2}{M^2}
=\log\frac{4k^2}{M^2}+O\!\left(\frac{M^2}{k^2}\right).
\end{equation}
As $M\to0$ at every fixed $k>0$, $r\to\infty$ and the symplectic
normal-mode transformation degenerates exactly as
$\operatorname{cond}(N)=e^{r}\,(1+o(1))\sim4k^2/M^2$ (condition number in
the Euclidean operator $2$-norm in the canonical basis); equivalently,
$\|\mu(S_+)\|\to\infty$: the positive representative escapes to infinity
in $\Sp(4,\C)/\Sp(2)$ (Section~\ref{sec:cartan}).  At the boundary the
dynamical map is a Jordan block and, by the Jordan obstruction of
\cite{companion1}, no positive-definite invariant form exists.
Precisely: the split-phase positive metric has \emph{no continuation as a
nondegenerate positive-definite invariant Hilbert metric} through
$c_1=0$.  This leaves open singular or positive-semidefinite boundary
forms, and the non-diagonalizable $PT$-symmetric Jordan realization of
Bender and Mannheim \cite{BM2008Jordan,Mannheim2021} remains available
\emph{as a distinct, singular completion of the boundary theory} --- it
is not a continuation of the split Hilbert metric, whose similarity
transformation itself becomes singular in the equal-frequency limit
\cite{BM2008Jordan}.  Only the Krein real form continues as a
nondegenerate completion of the same complex canonical variables; the
spectral functional continues distributionally on the curvature-generated
algebra (Definition~\ref{def:compmaps}).  (Numerically,
$\operatorname{cond}(N)=7.6,\,47.5,\,403,\,4447$ at
$M=1,\,0.3,\,0.1,\,0.03$, $k=1$, with
$\operatorname{cond}(N)/e^{r}\to1$; this serves as a regression check of
the exact law, not as the proof.)
\end{theorem}

\section{The master theorem}\label{sec:master}

\begin{theorem}[Classification]\label{thm:master}
Consider free scalar-free Einstein--Weyl gravity linearized about
Minkowski space.  Let $\mathfrak A_{\mathrm{inv}}$ be the algebra of
linear gauge-invariant observables, equivalently the reduced field
algebra, with the additional Weyl-null quotient imposed at $c_1=0$
(Definition~\ref{def:Ainv}).  Within the class of translation-invariant,
quasifree, mode-local constructions (Definition~\ref{def:realform})
compatible with the reduced symplectic form, the spectral condition, and
the Poincar\'e action on the physical polarization spaces, the following
hold.

For $c_1\neq0$, the complex spectral covariance is unique: the reduced
complex dynamical algebra $(\mathcal V_\C,\sigma,U)$ determines a unique
Poincar\'e-covariant spectral quasifree functional on
$\mathfrak A_{\mathrm{inv}}$.  Up to $U$-equivariant complex-symplectic
automorphisms preserving that covariance --- in particular, modulo
symplectic stabilizers commuting with the dynamics --- there are exactly
two admissible real-form classes (Definition~\ref{def:realform}):
\begin{enumerate}
\item a positive pseudo-Hermitian completion in which the full massive
spin-$2$ multiplet has uniformly rotated reality; and
\item a Krein completion in which the gravitational field retains its
standard reality and the massive spin-$2$ one-particle irrep carries
uniform negative sign relative to the positive massless graviton sector.
\end{enumerate}
No completion assigning different signatures or real structures to
different helicities of the massive spin-$2$ irrep is Poincar\'e
covariant.  The two admissible real forms induce the same complex
bilinear covariance, and hence the same complex correlation functions on
$\mathfrak A_{\mathrm{inv}}$; they differ only in involution and
completion (Remark~\ref{rem:states}).

As $c_1\to0$, the functionals pulled back to the curvature-generated
regular subalgebra converge distributionally,
$\omega_{c_1}\circ\iota_{c_1}\to\omega_0\circ\iota_0$
(Definition~\ref{def:compmaps}); no convergence is claimed on the
complete split-phase observable algebra, whose rank changes at the
boundary.  The transverse-traceless sector becomes a pair of $\Box^2$
Jordan systems, the vector sector becomes two ordinary second-order
massless ghost modes, and the scalar sector enters the kernel of the
presymplectic form and is removed by the emergent Weyl gauge symmetry.
The split-phase positive metric has no continuation as a nondegenerate
positive-definite invariant metric through this boundary (any
$PT$-symmetric Jordan realization there is a distinct, singular
completion), whereas the Krein real form continues as a nondegenerate
completion of the same complex canonical variables to the resulting
Jordan theory.
\end{theorem}

\begin{proof}
Category: Definition~\ref{def:realform}.  Uniqueness of the complex
covariance on $\mathfrak A_{\mathrm{inv}}$: Theorems~\ref{thm:kernel},
\ref{thm:reassembly} (residues forced by the reduced symplectic form and
covariance; the potential representative is fixed exactly modulo
pure-gauge kernels by \eqref{eq:exactkernel}).  Two real forms and no
hybrid: Theorem~\ref{thm:trilemma} (modewise), Lemma~\ref{lem:wigner}
(lift to the Poincar\'e group; two involution classes),
Theorems~\ref{thm:nohybrid}--\ref{thm:tworeal}.  Identity of complex
correlators: Theorems~\ref{thm:bridge4}, \ref{thm:weyl},
Remark~\ref{rem:states}.  Stabilizer statement: orbit constancy
\cite{companion3} with the enlarged stabilizers of
Section~\ref{sec:reduction} and the canonical representative of
Proposition~\ref{prop:cartan}.  Boundary: Theorems~\ref{thm:conformal},
\ref{thm:term}, Definition~\ref{def:compmaps}, and
Remark~\ref{rem:quotient}.
\end{proof}

\section{Discussion}\label{sec:discussion}

\paragraph{Relation to the two programs.}
The theorem places the positive-PT quantization \cite{BM2008PRD,
Mannheim2021} and the Krein quantization on the two sides of
a mathematically controlled classification: they are the two covariant
real forms of one reduced complex spectral theory, agreeing on every
complex gauge-invariant correlation function, and distinguished by which
member of the trilemma they sacrifice --- standard gravitational reality
(positive form) or positivity of the norm (Krein form).  A terminological
caution: Bateman and Turok \cite{BT2026} quantize the scalar
perfect-square theory, not this gravitational field; the second class is
therefore a \emph{Bateman--Turok-type Krein real form} --- the
gravitational extension of their construction --- rather than their
construction itself.  The conformal locus is one-sided in the precise
sense of Theorem~\ref{thm:term}: only the Krein real form continues as a
nondegenerate completion of the same complex canonical variables, while
the split-phase positive metric has no nondegenerate invariant
positive-definite limit, and any $PT$-symmetric Jordan realization at the
boundary \cite{BM2008Jordan} is a distinct, singular completion rather
than a continuation of the split Hilbert metric.  In this precise sense
Mannheim's construction is a split-phase real form and the
Bateman--Turok-type Krein form is its resonant continuation.  One
overreading must be excluded: agreement of the two real forms on every
complex gauge-invariant correlation function is a \emph{free-theory}
statement, and does not make the completions interchangeable
descriptions of an interacting theory.  For the fourth-order scalar the
companion interaction paper \cite{companionV} proves they separate:
the canonical positive metric is obstructed by on-shell branch-changing
conversion while the Krein/ghost-parity structure survives the same
interaction exactly.  The corresponding gravitational computation has
not been performed (see the outlook below).

\paragraph{Relation to Mannheim's CPT program.}
Mannheim has argued \cite{Mannheim2016} that time-independent
conservation of scalar products in a non-Hermitian relativistic theory
forces spectra that are real or occur in complex-conjugate pairs
together with an antilinear symmetry, which complex Lorentz invariance
identifies with CPT; the underlying equivalence of pseudo-Hermiticity,
antilinear symmetry, and real-or-conjugate-pair spectra (with matched
Jordan data) is Mostafazadeh's \cite{Mostafazadeh2002nd}.  None of that
organization is claimed here; the contribution of this paper is the
gauge-reduced helicity structure and its consequences.  Placed against
that program, the present classification says the following.  (i) Away
from the conformal locus every massive spin-$2$ polarization admits a
positive completion, but only \emph{jointly}: the trilemma and the
no-hybrid theorem force one uniform choice across the multiplet, and the
positive choice carries uniformly rotated massive reality
(Theorems~\ref{thm:trilemma}--\ref{thm:tworeal}).  A CPT-based positive
scalar product in the sense of \cite{Mannheim2016} therefore lies
\emph{inside} the classified positive real-form class --- with the
rotated-reality qualification made explicit here.  (ii) At the conformal
locus the positive representative escapes to infinity
($\operatorname{cond}(N)=e^{r}(1+o(1))$, Theorem~\ref{thm:term}) and the
TT dynamics becomes a Jordan block.  (iii) What the Jordan theorem
excludes there is exactly the class of nondegenerate, time-independent,
positive-definite invariant Hermitian forms making the Jordan generator
self-adjoint --- and nothing more.  The equal-frequency treatment of
Bender and Mannheim \cite{BM2008Jordan} uses nonstationary Jordan states
and a degenerate pairing, i.e.\ a construction \emph{outside} the
excluded class: there is no contradiction between the no-go and that
realization, only a sharp delimitation.  (iv) The trilemma is then the
precise form of the choice a CPT-based ghost-free quadratic gravity must
make away from the boundary: rotated massive reality with a positive
norm, or standard gravitational reality with a Krein norm --- no
covariant hybrid exists.

\paragraph{The geometric picture.}
The Cartan formulation of Section~\ref{sec:cartan} turns the phase
diagram into geometry: the split phase is the locus where the positive
representative sits at finite symmetric-space displacement
$\|\mu(S_+)\|=\sqrt2\,r$ from the origin, the approach to the conformal
locus is escape to infinity along the $C_2$ chamber wall, and the Jordan
boundary is the ideal boundary point at which no positive-definite
invariant form remains.  Minimum distortion, orbit constancy, and the
Jordan obstruction --- imported here from the companions --- are the
value, the fibers, and the boundary behaviour of the single map $\mu$
restricted to the split orbit.

\paragraph{Boundary truncation is a third, distinct construction.}
Maldacena's boundary-condition selection of the Einstein branch from
conformal gravity \cite{Maldacena} removes one branch of the solution
algebra, $\ker(L_1L_2)\to\ker L_1$; the completions classified here retain
the full fourth-order solution algebra and change its inner-product
structure.  The two mechanisms answer the ghost problem differently and
should not be identified; on Einstein backgrounds with $\Lambda\neq0$ the
critical locus shifts \cite{LuPope} and pure conformal gravity contains a
partially massless sector \cite{DeserWaldron}, so the flat-space phase
diagram above is the $\Lambda\to0$ slice of a richer picture we do not
analyze here.

\paragraph{Interaction diagnostic (outlook).}
The free classification fixes the arena for the interacting question: for
the positive real form, whether the rotated massive reality is compatible
with the cubic Weyl vertices; for the Krein form, whether a gravitational
ghost parity commutes with the interaction and with BRST symmetry, as its
scalar analogue does in the perfect-square theory \cite{BT2026}.  Both are
sharp, computable questions --- the natural first test is
$[\mathcal P_{\mathrm{ghost}},S_{\mathrm{int}}]=0$ at cubic order --- and
they provide the first necessary tests of whether either real form
supports an interacting unitarity statement; they do not by themselves
settle loop anomalies, BRST cohomology, renormalization of the parity, or
positivity of physical probabilities.  The companion interaction paper
\cite{companionV} supplies an interaction-obstruction mechanism ---
on-shell branch-changing conversion obstructing the deformation of the
positive metric --- that can now be applied to the cubic Einstein--Weyl
vertices; no interacting gravitational conclusion is assumed here.  We
have not performed these computations; the free theorem is neutral on
the outcome.

\begin{thebibliography}{12}

\bibitem{Stelle}
K.~S.~Stelle,
Phys.\ Rev.\ D \textbf{16} (1977) 953.

\bibitem{BM2008PRD}
C.~M.~Bender and P.~D.~Mannheim,
\emph{No-ghost theorem for the fourth-order derivative Pais--Uhlenbeck
oscillator model},
Phys.\ Rev.\ Lett.\ \textbf{100} (2008) 110402,
\href{https://doi.org/10.1103/physrevlett.100.110402}{doi:10.1103/physrevlett.100.110402}.

\bibitem{BM2008Jordan}
C.~M.~Bender and P.~D.~Mannheim,
\emph{Exactly solvable PT-symmetric Hamiltonian having no Hermitian
counterpart},
Phys.\ Rev.\ D \textbf{78} (2008) 025022, arXiv:0804.4190.

\bibitem{Mannheim2021}
P.~D.~Mannheim,
\emph{Solution to the ghost problem in higher-derivative gravity},
arXiv:2109.12743.

\bibitem{BT2026}
S.~Bateman and N.~Turok,
\emph{Escape from Ostrogradsky via hidden ghost parity},
arXiv:2607.00096.

\bibitem{companion1}
GPT-5.6.sol and A.~A.~Palm,
\emph{Canonical Positive Symplectic Diagonalization of the
Pais--Uhlenbeck Oscillator},
Paper~I of the pure-Weyl gravity programme, 2026;
file \path{paper/01-symplectic-diagonalization.pdf} in the repository
\url{https://github.com/area9innovation/weyl-gravity}.

\bibitem{companion2}
GPT-5.6.sol and A.~A.~Palm,
\emph{The Pais--Uhlenbeck Metric as a Minimum-Distortion Principle,
and the Representation Problem for the Fourth-Order Field},
Paper~II of the pure-Weyl gravity programme, 2026;
file \path{paper/02-variational-fock.pdf} in the repository
\url{https://github.com/area9innovation/weyl-gravity}.

\bibitem{companion3}
GPT-5.6.sol and A.~A.~Palm,
\emph{The Universal Vacuum of the Fourth-Order Scalar Field:
Metric Orbits, Fock Sectors, and the Krein Boundary},
Paper~III of the pure-Weyl gravity programme, 2026;
file \path{paper/03-fourth-order-vacuum.pdf} in the repository
\url{https://github.com/area9innovation/weyl-gravity}.

\bibitem{companionV}
GPT-5.6.sol and A.~A.~Palm,
\emph{Interaction Obstructions, Resonant PT Breaking,
and Doubled Jordan Symmetry in Fourth-Order Theories},
Paper~V of the pure-Weyl gravity programme, 2026;
file \path{paper/05-interaction-obstructions.pdf} in the repository
\url{https://github.com/area9innovation/weyl-gravity}.

\bibitem{Mannheim2016}
P.~D.~Mannheim,
\emph{CPT symmetry without Hermiticity},
arXiv:1611.02100.

\bibitem{Mostafazadeh2002nd}
A.~Mostafazadeh,
\emph{Pseudo-Hermiticity for a class of nondiagonalizable Hamiltonians},
J.\ Math.\ Phys.\ \textbf{43} (2002) 6343, arXiv:math-ph/0207009.

\bibitem{LuPope}
H.~L\"u and C.~N.~Pope,
\emph{Critical gravity in four dimensions},
Phys.\ Rev.\ Lett.\ \textbf{106} (2011) 181302, arXiv:1101.1971.

\bibitem{Maldacena}
J.~Maldacena,
\emph{Einstein gravity from conformal gravity},
arXiv:1105.5632.

\bibitem{DeserWaldron}
S.~Deser, E.~Joung and A.~Waldron,
\emph{Partial masslessness and conformal gravity},
J.\ Phys.\ A \textbf{46} (2013) 214019, arXiv:1208.1307.

\bibitem{Holdom}
B.~Holdom,
Nucl.\ Phys.\ B \textbf{1000} (2024) 116696.

\bibitem{Riegert}
R.~J.~Riegert,
Phys.\ Lett.\ B \textbf{134} (1984) 56.

\end{thebibliography}

\appendix

\section{Verification checks and the theorems they support}
\label{app:checks}

Scripts: \texttt{verify\_gravity\_reduction.py} (G1--G7),
\texttt{verify\_gravity\_completion.py} (G8--G9),
\texttt{verify\_gravity\_spectral.py} (G10--G12), each printing
\texttt{ALL PASS}; shared engine \texttt{gravity\_engine.py} (exact
$O(\epsilon^2)$ second variation).

\begin{center}
\begin{tabular}{l l}
\hline
Check & Supports\\
\hline
G1--G2 & TT perfect-square Lagrangian and PU block
  (Theorem~\ref{thm:tt})\\
G3 & vector $w$-reduction, unpaired massive ghost
  (Theorem~\ref{thm:lower}(i))\\
G4--G5 & scalar reduction, scalaron decoupling at $\alpha=-3\beta$
  (Theorem~\ref{thm:lower}(ii))\\
G6 & sector gauge invariance (Section~\ref{sec:reduction})\\
G7 & polarization-enhanced stabilizer $4\to16$
  (Remark after Definition~\ref{def:Ainv})\\
G8 & quarter-turn identities, trilemma rows, coset
  $T_0\cdot SO(2,\C)$ (Theorem~\ref{thm:trilemma})\\
G9a--b & $SO(3)$ commutant $=\R I_5$, hybrid violation
  (Theorem~\ref{thm:nohybrid})\\
G9c & $[J_i,D_{\mathrm{tot}}]=0$; quarter-turn is admissible positive
  diagonalizer (Theorem~\ref{thm:tworeal})\\
G10a--b & sector kernels, bisolution property, commutator
  normalizations (Theorem~\ref{thm:kernel})\\
G10c & projector contractions $\tfrac12:\tfrac{M^2}2:\tfrac16$;
  absolute coefficient $C=4/c_1$ (Theorem~\ref{thm:reassembly})\\
G10d & real-form independence of the complex kernel
  (Theorem~\ref{thm:bridge4})\\
G11a--c & Riemann kills pure gauge; Weyl traceless;
  $\Pi^{(2,M)}$--$\Pi_0$ contraction equality
  (Theorem~\ref{thm:weyl}(i)--(ii))\\
G11e & Weyl symbol nonzero on massless TT, kills null gauge
  (proof of Theorem~\ref{thm:weyl}(iv))\\
G12a, a$'$ & TT Jordan limit incl.\ the $1/\gamma$ prefactor
  (Theorem~\ref{thm:conformal}(i))\\
G12b & vector massless limit, fixed $\mu_V$
  (Theorem~\ref{thm:conformal}(ii))\\
G12c & Weyl shift is a symmetry iff $c_1=0$; scalar kinetic
  coefficient $\to0$ (Theorem~\ref{thm:conformal}(iii))\\
G12d, d$'$ & $\operatorname{cond}(N)$ divergence and exact law
  $\operatorname{cond}(N)/e^r\to1$ (Theorem~\ref{thm:term})\\
\hline
\end{tabular}
\end{center}

Checks establish identities; the representation-theoretic and microlocal
statements (Lemma~\ref{lem:wigner}, existence of completions,
$\WF$ equality) are analytic, as stated in the Verification paragraph.

\end{document}
